\section{Justificativa}
\label{sec:1.3}

A relevância prática deste trabalho reside na redução direta do esforço operacional imposto a coordenadores acadêmicos durante o processo de montagem semestral de grades. O tempo dedicado à resolução manual de conflitos de horários representa um custo institucional significativo, pois horas que poderiam ser direcionadas a atividades pedagógicas, acompanhamento de alunos e planejamento curricular acabam sendo consumidas por uma tarefa logística repetitiva e propensa a inconsistências. A automação desse processo, ainda que parcial e sujeita à validação humana posterior, representa um ganho concreto de eficiência operacional para as IES.

Do ponto de vista acadêmico, a contribuição do presente trabalho se apoia em dois eixos complementares. O primeiro é a aplicação de Algoritmos Genéticos ao UCTP em um sistema funcional completo, não apenas como protótipo algorítmico, mas integrado a módulos de gestão de entidades acadêmicas, controle de acesso por perfis, notificações e ciclo de vida de grades. Essa integração aproxima a solução proposta das condições reais de uso em uma IES. O segundo eixo é a comparação controlada de três implementações de Algoritmos Genéticos aplicadas ao mesmo problema, desenvolvidas em paradigmas distintos: orientado a objetos, com Java; funcional, com Python; e concorrente, com Go. Esse recorte é pouco explorado na literatura de timetabling, que em geral avalia o desempenho dos algoritmos sem aprofundar o impacto das escolhas de implementação \cite{premananda2024, bashab2023}.

A escolha da arquitetura de monolito modular com adaptadores de microsserviços para os algoritmos responde a uma necessidade técnica específica: isolar o processamento computacionalmente intensivo da geração de grades do restante do sistema, sem transformar toda a aplicação em uma arquitetura distribuída. Esse padrão permite que a API principal permaneça responsiva durante a execução dos Algoritmos Genéticos, que podem demandar diversas gerações evolutivas, além de viabilizar a substituição ou adição de novas implementações sem impacto direto nos demais módulos do sistema \cite{newman2021}.

Por fim, a adoção de um processo de desenvolvimento iterativo e incremental, com requisitos funcionais e não funcionais documentados e testados, confere ao IGRIS o caráter de artefato de engenharia de software, e não apenas de experimento computacional. A disponibilização do sistema e dos dados de benchmark contribui para que pesquisadores e instituições possam replicar, estender ou adaptar a solução a contextos acadêmicos específicos, ampliando o alcance prático e científico do trabalho \cite{sommerville2016}.
